\begin{textsample}{POS Dim 1 – persona\_grok – Score 17.00 – t451\_grok.txt}  \label{ex:f1_pos_036}
\textbf{Today} marks 15 \textbf{years} since the passing of Greenpeace co-founder Bob Hunter on May 2, 2005.
While best known for his role in establishing our \textbf{organization}, Bob 's life extended far beyond Greenpeace, encompassing his early \textbf{work} as a \textbf{writer} and \textbf{journalist}.
He drew profound influences from \textbf{ecology}, the \textbf{peace} \textbf{movements}, and \textbf{media} theorists such as Marshall McLuhan.
At the heart of Bob 's approach was a philosophy of non-violent \textbf{revolution}, achieved through what he called"\textbf{mind} bombs"—powerful \textbf{images} designed to elevate ecological consciousness.
This \textbf{vision} drove key \textbf{campaigns}, including protests against nuclear \textbf{tests} and efforts to save \textbf{whales}.
Bob often spoke of spirituality and the interconnectedness of nature, viewing \textbf{ecology} as a subversive science capable of challenging established norms.
He openly \textbf{shared} his personal struggles, and after retiring from Greenpeace in 1977, he continued his contributions through writing \textbf{books} and receiving awards that recognized his enduring impact.
His legacy remains a guiding force in our ongoing \textbf{work} for the planet.

% matched lemmas: book, campaign, ecology, image, journalist, medium, mind, movement, organization, peace, revolution, share, test, today, vision, whale, work, writer, year
\end{textsample}
